\documentclass{article}
\usepackage[utf8]{inputenc}
\usepackage{a4wide}

\begin{document}

\section*{ย่องเบาหลบกับระเบิด (mine)}

\section{\texorpdfstring{\textbf{ย่องเบาหลบกับระเบิด!}}{ย่องเบาหลบกับระเบิด!}}\label{uxe22uxe2duxe07uxe40uxe1auxe32uxe2buxe25uxe1auxe01uxe1auxe23uxe30uxe40uxe1auxe14}



\subparagraph{\texorpdfstring{กำหนดสนาม 4 เหลี่ยมมีพื้นที่RxCมีกับระเบิดอยู่จำนวนหนึ่ง

ให้หาทางเดินจากด้านซ้ายสุดไปด้านขวาสุดโดย\ul{{\textbf{ไม่ไปอยู่ข้างๆ}}}กับระเบิดนั้นทั้งนี้การเดินจะเดินได้เพียง4ทิศ

คือ ซ้ายขวาบนหรือล่างไม่สามารถเดินทแยงได้ หากเดินไปข้างๆ รอบ 8

ทิศของกับระเบิดก็จะตาย\textbf{โจทย์}คือให้หาจำนวน{\textbf{ก้าว}}เดินของเส้นทางที่ปลอดภัยและ\textbf{สั้นที่สุด}ที่ต้องเดินจากขอบซ้ายสุด(แถวแนวนอนแถวใดก็ได้)

ไปจนถึงขอบขวาสุด (จะแถวใดก็ได้)หากไม่มีเส้นทางเดินที่ปลอดภัย ให้ตอบว่า\textbf{No

route}เช่นในตัวอย่างที่ 4 (เพราะกับระเบิดมากไปนิด

เลยไม่เหลือทางปลอดภัยให้เดิน)}{กำหนดสนาม 4 เหลี่ยมมีพื้นที่RxCมีกับระเบิดอยู่จำนวนหนึ่ง ให้หาทางเดินจากด้านซ้ายสุดไปด้านขวาสุดโดยไม่ไปอยู่ข้างๆกับระเบิดนั้นทั้งนี้การเดินจะเดินได้เพียง4ทิศ คือ ซ้ายขวาบนหรือล่างไม่สามารถเดินทแยงได้ หากเดินไปข้างๆ รอบ 8 ทิศของกับระเบิดก็จะตายโจทย์คือให้หาจำนวนก้าวเดินของเส้นทางที่ปลอดภัยและสั้นที่สุดที่ต้องเดินจากขอบซ้ายสุด(แถวแนวนอนแถวใดก็ได้) ไปจนถึงขอบขวาสุด (จะแถวใดก็ได้)หากไม่มีเส้นทางเดินที่ปลอดภัย ให้ตอบว่าNo routeเช่นในตัวอย่างที่ 4 (เพราะกับระเบิดมากไปนิด เลยไม่เหลือทางปลอดภัยให้เดิน)}}\label{uxe01uxe33uxe2buxe19uxe14uxe2auxe19uxe32uxe21-4-uxe40uxe2buxe25uxe22uxe21uxe21uxe1euxe19uxe17rxcuxe21uxe01uxe1auxe23uxe30uxe40uxe1auxe14uxe2duxe22uxe08uxe33uxe19uxe27uxe19uxe2buxe19uxe07-uxe43uxe2buxe2buxe32uxe17uxe32uxe07uxe40uxe14uxe19uxe08uxe32uxe01uxe14uxe32uxe19uxe0buxe32uxe22uxe2auxe14uxe44uxe1buxe14uxe32uxe19uxe02uxe27uxe32uxe2auxe14uxe42uxe14uxe22uxe44uxe21uxe44uxe1buxe2duxe22uxe02uxe32uxe07uxe46uxe01uxe1auxe23uxe30uxe40uxe1auxe14uxe19uxe19uxe17uxe07uxe19uxe01uxe32uxe23uxe40uxe14uxe19uxe08uxe30uxe40uxe14uxe19uxe44uxe14uxe40uxe1euxe22uxe074uxe17uxe28-uxe04uxe2d-uxe0buxe32uxe22uxe02uxe27uxe32uxe1auxe19uxe2buxe23uxe2duxe25uxe32uxe07uxe44uxe21uxe2auxe32uxe21uxe32uxe23uxe16uxe40uxe14uxe19uxe17uxe41uxe22uxe07uxe44uxe14-uxe2buxe32uxe01uxe40uxe14uxe19uxe44uxe1buxe02uxe32uxe07uxe46-uxe23uxe2duxe1a-8-uxe17uxe28uxe02uxe2duxe07uxe01uxe1auxe23uxe30uxe40uxe1auxe14uxe01uxe08uxe30uxe15uxe32uxe22uxe42uxe08uxe17uxe22uxe04uxe2duxe43uxe2buxe2buxe32uxe08uxe33uxe19uxe27uxe19uxe01uxe32uxe27uxe40uxe14uxe19uxe02uxe2duxe07uxe40uxe2auxe19uxe17uxe32uxe07uxe17uxe1buxe25uxe2duxe14uxe20uxe22uxe41uxe25uxe30uxe2auxe19uxe17uxe2auxe14uxe17uxe15uxe2duxe07uxe40uxe14uxe19uxe08uxe32uxe01uxe02uxe2duxe1auxe0buxe32uxe22uxe2auxe14uxe41uxe16uxe27uxe41uxe19uxe27uxe19uxe2duxe19uxe41uxe16uxe27uxe43uxe14uxe01uxe44uxe14-uxe44uxe1buxe08uxe19uxe16uxe07uxe02uxe2duxe1auxe02uxe27uxe32uxe2auxe14-uxe08uxe30uxe41uxe16uxe27uxe43uxe14uxe01uxe44uxe14uxe2buxe32uxe01uxe44uxe21uxe21uxe40uxe2auxe19uxe17uxe32uxe07uxe40uxe14uxe19uxe17uxe1buxe25uxe2duxe14uxe20uxe22-uxe43uxe2buxe15uxe2duxe1auxe27uxe32no-routeuxe40uxe0auxe19uxe43uxe19uxe15uxe27uxe2duxe22uxe32uxe07uxe17-4-uxe40uxe1euxe23uxe32uxe30uxe01uxe1auxe23uxe30uxe40uxe1auxe14uxe21uxe32uxe01uxe44uxe1buxe19uxe14-uxe40uxe25uxe22uxe44uxe21uxe40uxe2buxe25uxe2duxe17uxe32uxe07uxe1buxe25uxe2duxe14uxe20uxe22uxe43uxe2buxe40uxe14uxe19}



\subparagraph{กำหนดค่าRและC∈~{[}10...900{]}

และจำนวนระเบิด∈{[}2...50{]}}\label{uxe01uxe33uxe2buxe19uxe14uxe04uxe32ruxe41uxe25uxe30c-10...900-uxe41uxe25uxe30uxe08uxe33uxe19uxe27uxe19uxe23uxe30uxe40uxe1auxe142...50}



\textbf{{ตัวอย่าง}}



หากกำหนดแผนผังสนามดังรูปซ้ายมือ. ให้ 0 เป็นกับระเบิดดังนั้นตำแหน่งที่เดินได้จะเป็น 1

ตามรูปขวามือ\ul{สังเกต}ว่าแต่ละช่อง 0 ในรูปซ้ายมือจะขยายกว้างออก 1

ช่องโดยรอบดังในรูปขวามือ.



เส้นทางสีเขียวคือคำตอบ และตอบ\textbf{11 (ไม่นับตำแหน่งแรกคอลัมน์ซ้ายมือสุด)}



\includegraphics[width=5.72917in,height=2.22917in,alt={image.png}]{/public/upload/e353966134.png}\\



{\hfill\break

}



{\textbf{ตัวอย่างที่ 4{ตอบ}{13}}}



{}



\includegraphics[width=2.08333in,height=4.22917in,alt={image.png}]{/public/upload/9166f1d0d9.png}\\



\hfill\break



\textbf{ข้อกำหนด}



{\def\LTcaptype{none} % do not increment counter

\begin{longtable}[]{@{}

  >{\raggedright\arraybackslash}p{(\linewidth - 2\tabcolsep) * \real{0.5000}}

  >{\raggedright\arraybackslash}p{(\linewidth - 2\tabcolsep) * \real{0.5000}}@{}}

\toprule\noalign{}

\begin{minipage}[b]{\linewidth}\centering

หัวข้อ

\end{minipage} & \begin{minipage}[b]{\linewidth}\centering

เงื่อนไข

\end{minipage} \\

\midrule\noalign{}

\endhead

\bottomrule\noalign{}

\endlastfoot

\begin{minipage}[t]{\linewidth}\raggedright

ข้อมูลนําเข้า\\

\strut

\end{minipage} & \begin{minipage}[t]{\linewidth}\raggedright

Standard Input (คีย์บอร์ด)\\

\strut

\end{minipage} \\

\begin{minipage}[t]{\linewidth}\raggedright

ข้อมูลส่งออก\\

\strut

\end{minipage} & \begin{minipage}[t]{\linewidth}\raggedright

Standard Output (จอภาพ)\\

\strut

\end{minipage} \\

\begin{minipage}[t]{\linewidth}\raggedright

ระยะเวลาสูงสุดที่ใช้ในการประมวลผล\\

\strut

\end{minipage} & 1000 ms \\

\begin{minipage}[t]{\linewidth}\raggedright

หน่วยความจําสูงสุดที่ใช้ในการประมวลผล\\

\strut

\end{minipage} & 64 MB \\

คะแนนสูงสุดของโจทย์ & 100 \\

\end{longtable}

}



\subsection*{Input}
{\def\LTcaptype{none} % do not increment counter

\begin{longtable}[]{@{}

  >{\raggedright\arraybackslash}p{(\linewidth - 2\tabcolsep) * \real{0.5000}}

  >{\raggedright\arraybackslash}p{(\linewidth - 2\tabcolsep) * \real{0.5000}}@{}}

\toprule\noalign{}

\begin{minipage}[b]{\linewidth}\raggedright

ข้อมูลนำเข้า

\end{minipage} & \begin{minipage}[b]{\linewidth}\raggedright

\hfill\break

\strut

\end{minipage} \\

\midrule\noalign{}

\endhead

\bottomrule\noalign{}

\endlastfoot

บรรทัดที่ 1 & \begin{minipage}[t]{\linewidth}\raggedright

มีจำนวนเต็มสองค่ากำหนดค่า R และ C ตามลำดับ คั่นด้วยช่องว่าง\\

\strut

\end{minipage} \\

บรรทัดที่ 2 ถึง R + 1 & \begin{minipage}[t]{\linewidth}\raggedright

เป็นค่าของพื้นสนามแต่ละช่อง คั่นด้วยช่องว่าง โดย 0เป็น ระเบิด และ 1 หมายถึงไม่ใช่ระเบิด\\

\strut

\end{minipage} \\

\end{longtable}

}



\subsection*{Output}
{\def\LTcaptype{none} % do not increment counter

\begin{longtable}[]{@{}

  >{\raggedright\arraybackslash}p{(\linewidth - 2\tabcolsep) * \real{0.5000}}

  >{\raggedright\arraybackslash}p{(\linewidth - 2\tabcolsep) * \real{0.5000}}@{}}

\toprule\noalign{}

\begin{minipage}[b]{\linewidth}\raggedright

ข้อมูลส่งออก

\end{minipage} & \begin{minipage}[b]{\linewidth}\raggedright

\hfill\break

\strut

\end{minipage} \\

\midrule\noalign{}

\endhead

\bottomrule\noalign{}

\endlastfoot

บรรทัดที่ 1 & \begin{minipage}[t]{\linewidth}\raggedright

จำนวนเต็มแสดงจำนวนก้าวขั้นต่ำที่จะต้องเดินจากขอบซ้ายไปขอบขวา และตอบว่า No route

หากไม่มีเส้นทางเดินที่ปลอดภัย\\

\strut

\end{minipage} \\

\end{longtable}

}



\end{document}
